%% ===========================================================================
%%  Dirección pública del demo desplegado — contrato versionado
%%
%%  Las URL de Cloud Run llevan el identificador del proyecto de GCP y cambian
%%  si el servicio se recrea, así que por decisión de E0 no se versionan. Este
%%  archivo es la mitad que SÍ vive en el repositorio: define siempre las tres
%%  secuencias de control, exista o no la dirección. El valor real vive en
%%  datos/despliegue.tex, que .gitignore excluye y que se escribe localmente
%%  antes de compilar la entrega:
%%
%%      bash scripts/escribir_url_demo.sh
%%
%%  Sin ese archivo el documento compila igual y dice, en lugar del enlace,
%%  que la dirección se entrega junto con el documento. Ningún .tex versionado
%%  queda roto por su ausencia, que es la razón de que el condicional se
%%  resuelva aquí una sola vez y no dentro de cada macro: un \if suelto dentro
%%  de una celda de tabularx rompe la alineación.
%%
%%  Cargarlo en el preámbulo, después de \usepackage{estilo/uxdoc} porque
%%  \url viene de hyperref. Compilar desde docs/entregables/.
%% ===========================================================================
\newif\ifhaydemoweb
\haydemowebfalse
\newcommand{\urlDemoWebValor}{}

\IfFileExists{datos/despliegue.tex}{%
  \input{datos/despliegue.tex}%
  \haydemowebtrue
}{}

%% El condicional se decide AQUI, en el preámbulo, y las dos macros salen ya
%% resueltas, de modo que en el punto de uso solo queda texto. La alternativa
%% —una sola definición con el \ifhaydemoweb dentro del cuerpo— dejaría un
%% condicional vivo en cada punto de uso, incluidas las celdas de tabularx, que
%% reexpande su cuerpo para medir las columnas. No hace falta averiguar si ese
%% caso concreto rompe: resolverlo una vez aquí es igual de corto y no deja la
%% pregunta abierta.
%%
%% Comprobado el 22-ago-2026 en los dos escenarios, con cero errores de TeX:
%% con datos/despliegue.tex presente el documento imprime el enlace y toma la
%% rama «con dirección»; sin el archivo imprime la frase de respaldo, toma la
%% rama «sin dirección» y el PDF resultante no contiene ninguna URL.
%%
%% \siHayDemoWeb{texto con dirección}{texto sin dirección}
%%   Elige entre dos redacciones completas. Úsese cuando la frase cambia de
%%   forma, no solo de dato.
%% \urlDemoWeb
%%   La dirección como enlace, o el respaldo en prosa. NO usarla dentro de
%%   \caption ni de ningún otro argumento móvil: \url no sobrevive ahí.
\ifhaydemoweb
  \newcommand{\siHayDemoWeb}[2]{#1}
  \newcommand{\urlDemoWeb}{\url{\urlDemoWebValor}}
\else
  \newcommand{\siHayDemoWeb}[2]{#2}
  \newcommand{\urlDemoWeb}{la dirección pública se entrega junto con este documento}
\fi
